\subsection{Conductor equalities in odd characteristic}
\label{subsec:odd-higher-conductor}

The canonical divisor in the structure theorem is divisible by the
conductor minus one as a Cartier class. Adjunction on an Albanese
fibre combines this divisibility with Tate's genus-change theorem
and the bounds of Gu, Sun and Zhou to leave at most three possible
fibre genera.

\begin{proposition}\label{prop:odd-higher-conductor-equality}
Let $X$ be a smooth projective integral surface over an algebraically
closed field $k$ of odd characteristic $p$, let $Y\subset X$ be a smooth
integral very ample curve, and let $M$ be a nontrivial line bundle of
order $p$ whose restriction to $Y$ is trivial. Let $\omega$ be its
logarithmic form, let $pa$ be its exact vanishing order along $Y$, and
let $\pi:N\to X$ be the normalization of its $\mu_p$-torsor. Suppose
that
\[
 2g(Y)-2=paY^2.
\]
Put $d=Y^2$. Then $a\geq5$, $N$ is smooth with ample
canonical bundle, and
\begin{equation}\label{eq:odd-conductor-invariants}
\begin{gathered}
 K_X\sim(pa-1)Y,\qquad K_N\sim(a-1)\pi^*Y,\qquad p\mid d,\\
 K_N^2=p(a-1)^2d,\qquad c_2(N)=-pa\,d,\\
 12\chi(\OO_N)=p(a^2-3a+1)d.
\end{gathered}
\end{equation}
Let $f:N\to C$ be its Albanese fibration, let $g$ be its
arithmetic fibre genus, and put
\[
 L=\pi^*Y,\qquad \ell=L\cdot F
\]
for a general fibre $F$. Exactly the following three rows can
survive the canonical divisibility, genus-change and slope bounds:
\begin{equation}\label{eq:odd-conductor-three-genera}
\begin{array}{c|c|c|c}
 \ell&\text{range of }a&g&\text{divisibility}\\ \hline
 1&a\geq5\text{ odd}&(a+1)/2&p-1\mid a+1\\
 2&a\geq7&a&p-1\mid2a\\
 3&a\geq21\text{ odd}&(3a-1)/2&p-1\mid3a-1
\end{array}
\end{equation}
The second and third rows require a nonhyperelliptic fibration.
In particular $a\ne6$, and $a=5$ forces $g=3$ and
$p\in\{3,7\}$. Every case satisfies
\begin{equation}\label{eq:odd-conductor-characteristic-bound}
 p\leq2a+1.
\end{equation}
The rows describe necessary numerical conditions and do not assert
the existence of equality cases.
\end{proposition}

\begin{proof}
Theorem~\ref{thm:conductor-equality-structure} gives the identities
\eqref{eq:odd-conductor-invariants}, the smoothness of $N$, and its
ample canonical bundle. Thus $N$ is minimal of general type,
$c_2(N)<0$ and $\chi(\OO_N)>0$.
The negative-Chern-number theorem
\cite[Theorems~2 and~7 and their proofs, pp.~264--266]{ShepherdBarronGeography}
gives an Albanese fibration $f:N\to C$, with $C$ smooth,
$f_*\OO_N=\OO_C$, arithmetic fibre genus $g\geq2$, and
integral singular rational geometric generic fibre.

The fibre need not be smooth. A general closed fibre $F$ is a
Cartier curve, with $F^2=0$, and adjunction gives
\[
 (a-1)L\cdot F=K_N\cdot F=2g-2.
\]
The dualizing degree formula applies to a Gorenstein curve, so
singular fibres cause no change here
\cite[Tags~0B4A, 0C19 and~0BS6]{Stacks}.
The morphism to the smooth curve $C$ is flat: its local rings
are torsion free over the base discrete valuation rings.
Consequently its arithmetic fibre genus equals its generic
arithmetic genus. The generic curve is regular by localization,
and is proper and geometrically integral over $k(C)$.
Its geometric normalization is rational. Tate's genus-change
theorem, in the form of
\cite[Corollary~2.3, author p.~4]{SchroerGenusChange}, gives
\[
 p-1\mid2g.
\]

Set $s=a-1\geq4$. Since $L$ is an ample Cartier divisor,
\begin{equation}\label{eq:odd-conductor-integer-fibre-degree}
 \ell=L\cdot F\in\mathbb Z_{>0},\qquad
 g=1+\frac{s\ell}{2}.
\end{equation}
In particular $g=2$ is impossible. The exact Chern ratio is
\[
 R(a)=\frac{K_N^2}{\chi(\OO_N)}
      =\frac{12s^2}{s^2-s-1}.
\]
For a hyperelliptic fibration, Gu--Sun--Zhou's bound is
\[
 R(a)\leq h(g):=
 \frac{(12g+8)(g-1)}{g^2-g-1}.
\]
For a nonhyperelliptic fibration it is $18$ at $g=3$,
$840/47$ at $g=4$, and
\[
 R(a)\leq n(g):=
 \frac{12(g-1)(3g^2-4g-4)}{g(3g^2-12g+15)}
 \qquad(g\geq5)
\]
\cite[Theorem~3.2 and its proof, pp.~622--623]{GuSunZhou}.
All denominators below are positive in their stated ranges.

We first record the monotonicity and comparison needed to make
the genus range finite:
\[
 h'(g)=-\frac{4(2g^2+2g+1)}{(g^2-g-1)^2}<0
 \qquad(g\geq2),
\]
\[
 n'(g)=-\frac{4\{5g^2((g-3)^2-2)+4(g-1)(3g-5)\}}
 {(g^3-4g^2+5g)^2}<0
 \qquad(g\geq5),
\]
\[
 n(g)-h(g)=
 \frac{4(g-1)(3g+2)(g^2-4g+2)}
 {g(g^2-4g+5)(g^2-g-1)}>0
 \qquad(g\geq5).
\]
In the derivative of $n$, both summands in braces are positive
for $g\geq5$.

Define
\[
 P(s,g)=(s^2-s-1)(5g^2-15g+4)
       -(s+1)g(3g^2-12g+15).
\]
Direct subtraction gives
\begin{equation}\label{eq:odd-conductor-slope-polynomial}
 n(g)-R(a)=
 \frac{12P(s,g)}{(s^2-s-1)g(3g^2-12g+15)}.
\end{equation}
If $\ell\geq4$, then $g\geq2s+1\geq9$.
Both slope bounds are at most $n(2s+1)$, whereas
\[
 P(s,2s+1)=-2s(2s^3+21s^2+2s-5)<0.
\]
Equation~\eqref{eq:odd-conductor-slope-polynomial} gives a
contradiction. Hence $\ell\leq3$.

If $\ell=1$, then $s$ is even and $g=1+s/2=(a+1)/2$.
For every even $s\geq4$,
\[
 h(1+s/2)-R(a)=
 \frac{4s(s^2-s-10)}{(s^2-s-1)(s^2+2s-4)}>0.
\]
The nonhyperelliptic bound also exceeds $h(g)$: for $g=3,4$
this follows from $18>88/5$ and $840/47>168/11$;
for $g\geq5$ it follows from the displayed comparison of
$n$ and $h$. Thus neither slope envelope further excludes
this row. Tate's condition becomes $p-1\mid a+1$.

For $\ell\geq2$ one has $g\geq s+1$, and
\[
 h(s+1)-R(a)=
 -\frac{4s(s^2+5s+5)}{(s^2-s-1)(s^2+s-1)}<0.
\]
Monotonicity of $h$ therefore excludes the hyperelliptic case
in the remaining rows.

If $\ell=2$, then $g=s+1\geq5$ and
\[
 P(s,s+1)=s(2s^3-10s^2-3s+5).
\]
The cubic factor is $-39$ at $s=4$ and $-10$ at $s=5$.
For $s\geq6$ it is
\[
 2s^2(s-5)-3s+5\geq2s^2-3s+5>0.
\]
The slope inequality consequently holds exactly for $a\geq7$,
and Tate's condition is $p-1\mid2a$.

If $\ell=3$, then $s$ is even, $g=1+3s/2\geq7$, and
\[
 P(s,1+3s/2)=\frac{3s}{8}(3s^3-59s^2-8s+20).
\]
For $4\leq s\leq19$ the cubic factor is at most
$-2s^2-8s+20<0$, because $3s-59\leq-2$.
For $s\geq20$ it is at least
$s^2-8s+20=(s-4)^2+4>0$.
Thus this row is permitted exactly for odd $a\geq21$,
and its Tate condition is $p-1\mid3a-1$.
This proves the complete numerical reduction
\eqref{eq:odd-conductor-three-genera}.

In the first row $p\leq a+2$, and in the second
$p\leq2a+1$. In the third, $p-1$ is a proper divisor
of $3a-1$: equality would give the composite number $p=3a$.
It follows that
\[
 p\leq\frac{3a+1}{2}<2a+1.
\]
This proves \eqref{eq:odd-conductor-characteristic-bound}.
When $a=5$ only the first row remains, giving $g=3$ and
$p-1\mid6$, hence $p=3$ or $7$. None of the rows allows $a=6$.
\end{proof}

\begin{corollary}\label{cor:odd-conductor-numerical-threshold}
In Proposition~\ref{prop:odd-higher-conductor-equality}, the following
are necessary lower bounds on the conductor:
\[
 a\geq A_{\mathrm{num}}(p):=
 \begin{cases}
 5,&p=3\text{ or }7,\\
 7,&p=5,\\
 9,&p=11,\\
 11,&p=13,\\
 (p-1)/2,&p\geq17.
 \end{cases}
\]
Each right-hand side is the first integer for which some genus
and fibration type satisfy the numerical tests in the proposition.
For $p\geq17$, equality in this numerical lower bound forces
$g=a=(p-1)/2$ and a nonhyperelliptic Albanese fibration.
For example,
\[
\begin{array}{c|rrrrrrrrrr}
 p&3&5&7&11&13&17&19&23&29&31\\ \hline
 A_{\mathrm{num}}(p)&5&7&5&9&11&8&9&11&14&15\\
 g\text{ at this value}&3&4&3&5&6&8&9&11&14&15
\end{array}
\]
These thresholds do not assert the existence of geometric equality
cases.
\end{corollary}

\begin{proof}
Put $q=(p-1)/2$. In the first row of
\eqref{eq:odd-conductor-three-genera}, $a$ is odd and
$a\equiv-1\pmod{p-1}$. In the second, $a$ is a multiple of
$q$ and $a\geq7$. The third starts at $a=21$.
These conditions give the first five cases in the statement
directly.

If $p\geq17$, then $q\geq8$. The first row has
$a\geq p-2=2q-1\geq q$, and the second has $a\geq q$.
If the third had $a<q$, the positive integer
$(3a-1)/(2q)$ would be less than $3/2$, hence equal to one.
This would imply $p=3a$, impossible for the prime $p$.
Thus every row requires $a\geq q$.
The pair $a=g=q$ satisfies the second row and its slope bound,
as proved above. The first row cannot have $a=q$ for $q\geq8$,
and the third would require $(3q-1)/(2q)$ to be an integer,
which it is not. This proves all assertions.
\end{proof}

For any fixed conductor, the complete characteristic test is the
union of the three divisibility conditions in
\eqref{eq:odd-conductor-three-genera}, with their stated ranges.
It can have gaps above the first permitted conductor. The third
row is impossible when $p\equiv1\pmod3$.
All assertions in this subsection concern equality with the exact
conductor; they impose no replacement inequality on cases with
strict excess in $2g(Y)-2\geq paY^2$.
